\documentclass[11pt]{article}
\usepackage[margin=18mm,headheight=22pt]{geometry}
\usepackage{amsmath}
\usepackage{fontspec}
\usepackage{unicode-math}
\setmainfont{NewCM10}[Extension=.otf,UprightFont=*-Book,BoldFont=*-Bold,ItalicFont=*-BookItalic,BoldItalicFont=*-BoldItalic]
\setsansfont{NewCMSans10}[Extension=.otf,UprightFont=*-Book,BoldFont=*-Bold,ItalicFont=*-BookOblique,BoldItalicFont=*-BoldOblique]
\setmathfont{NewCMMath-Book.otf}
\usepackage{xcolor}
\definecolor{FAUblue}{HTML}{04316A}
\definecolor{Paperblue}{HTML}{EEF2F7}
\definecolor{EEIgold}{HTML}{CC9C42}
\usepackage{etoolbox,enumitem,booktabs,tabularx,array,needspace,microtype}
\usepackage[most]{tcolorbox}
\usepackage{titlesec,fancyhdr}
\usepackage[hidelinks]{hyperref}
\usepackage{draftwatermark}
\SetWatermarkText{Made by Thanate Amornjiraporn (Focus)}
\SetWatermarkColor{black!7}\SetWatermarkFontSize{30pt}\SetWatermarkAngle{45}
\setlength{\parindent}{0pt}
\setlength{\parskip}{5pt}
\setlist{itemsep=2pt,topsep=3pt,leftmargin=1.6em,beginpenalty=10000,endpenalty=10000}
\raggedbottom
\clubpenalty=10000 \widowpenalty=10000 \displaywidowpenalty=10000 \interlinepenalty=10000
\setcounter{secnumdepth}{0}
\setcounter{tocdepth}{1}
\pagestyle{fancy}\fancyhf{}
\fancyhead[L]{\small\sffamily\color{FAUblue} From Fundamentals to AI Apocalypse · APDR · Blackboard Exercise 3}
\fancyhead[R]{\small\sffamily Blackboard Exercise 3: Worked solutions}
\fancyfoot[C]{\small\thepage}
\titleformat{\section}{\Large\bfseries\sffamily\color{FAUblue}}{}{0pt}{}
\titleformat{\subsection}{\large\bfseries\sffamily\color{FAUblue}}{}{0pt}{}
\titlespacing*{\section}{0pt}{16pt}{6pt}
\preto\section{\Needspace{8\baselineskip}}
\preto\subsection{\Needspace{6\baselineskip}}
\newtcolorbox{equationcard}[1]{colback=Paperblue,colframe=Paperblue,boxrule=0pt,arc=0pt,left=10pt,right=10pt,top=6pt,bottom=6pt,title={#1},fonttitle=\small\bfseries\sffamily,coltitle=FAUblue,before skip=8pt,after skip=8pt}
\newtcolorbox{workedbox}[1]{colback=white,colframe=EEIgold,boxrule=0pt,leftrule=3pt,arc=0pt,left=10pt,top=4pt,bottom=4pt,title={#1},fonttitle=\small\bfseries\sffamily,coltitle=FAUblue,attach title to upper={\par},before skip=8pt,after skip=8pt}
\newtcolorbox{recallbox}[1]{colback=white,colframe=FAUblue,boxrule=0pt,leftrule=3pt,arc=0pt,left=10pt,top=4pt,bottom=4pt,title={#1},fonttitle=\small\bfseries\sffamily,coltitle=FAUblue,attach title to upper={\par},before skip=8pt,after skip=8pt}
\usepackage{environ}
\newsavebox\measurebox
% Typeset the section once invisibly to measure it; start a new page if it will not fit in what is left.
\makeatletter
\newif\ifkeep@measuring
\NewEnviron{keepsection}{%
  \ifkeep@measuring \BODY \else
    \edef\keep@sec{\the\value{section}}\edef\keep@sub{\the\value{subsection}}%
    \setbox\measurebox\vbox{\hsize\linewidth\keep@measuringtrue\def\Needspace##1{}\BODY}%
    \setcounter{section}{\keep@sec}\setcounter{subsection}{\keep@sub}%
    \dimen2=\ht\measurebox \advance\dimen2\dp\measurebox
    \dimen@=\textheight
    \ifdim\pagegoal<\maxdimen
      \dimen@=\pagegoal \advance\dimen@-\pagetotal
      \ifdim\dimen2>\dimen@ \clearpage \dimen@=\textheight \fi
    \fi
    % if it fits, nothing inside may force an early break
    \ifdim\dimen2>\dimen@ \BODY \else \begingroup\def\Needspace##1{}\BODY\endgroup \fi
  \fi}
\makeatother
\newtcolorbox{notebox}{colback=Paperblue,colframe=Paperblue,boxrule=0pt,arc=0pt,left=8pt,right=8pt,top=4pt,bottom=4pt,before skip=6pt,after skip=6pt}
\newtcolorbox{plainbox}{colback=white,colframe=FAUblue,boxrule=0pt,leftrule=3pt,arc=0pt,left=10pt,right=8pt,top=4pt,bottom=4pt,before skip=6pt,after skip=8pt}
\usetikzlibrary{arrows.meta}
\title{\sffamily\bfseries\color{FAUblue}Blackboard Exercise 3: Worked solutions}
\author{Algorithms, Programming, and Data Representation \ensuremath{\cdot} Module 92313 \ensuremath{\cdot} Blackboard Exercise 3\\[2pt]\small From Fundamentals to AI Apocalypse · personal study guide}
\date{}
\begin{document}
\maketitle
\begingroup\footnotesize\setlength{\parskip}{0pt}\tableofcontents\endgroup
\clearpage

\Needspace{8\baselineskip}
{\large Work through memory size, addressing, module construction, interleaving and byte order. For address calculations, state whether an address selects a word or a byte.}


\begin{keepsection}
\Needspace{5\baselineskip}
\section{Exercise 3.1: Memory chip architecture}

\textbf{Connected topic:} Addressing and capacity (§5.2)

\begin{keepsection}
\Needspace{6\baselineskip}
\subsection{3.1.1 Memory size in bits}

\Needspace{3\baselineskip}
Try first: For a 4G \ensuremath{\times} 32-bit memory, find the total capacity in bits and write it as 2\textsuperscript{x}.

\nopagebreak[4]{\small\itshape Answers: see “Exercise 3.1: Memory chip architecture” under Answers at the end.}
\end{keepsection}

\begin{keepsection}
\Needspace{6\baselineskip}
\subsection{3.1.2 Word-addressable memory}

\Needspace{3\baselineskip}
Try first: Find the address width when each address selects one 32-bit word.

\nopagebreak[4]{\small\itshape Answers: see “Exercise 3.1: Memory chip architecture” under Answers at the end.}
\end{keepsection}

\begin{keepsection}
\Needspace{6\baselineskip}
\subsection{3.1.3 Byte-addressable memory}

\Needspace{3\baselineskip}
Try first: Find the address width when each address selects one byte.

\nopagebreak[4]{\small\itshape Answers: see “Exercise 3.1: Memory chip architecture” under Answers at the end.}
\end{keepsection}

\begin{keepsection}
\Needspace{7\baselineskip}
\subsection{3.1.4 Byte address range}

\Needspace{4\baselineskip}
Try first: Give the lowest and highest byte addresses in hexadecimal and the number of bits needed to represent them.

\nopagebreak[4]{\small\itshape Answers: see “Exercise 3.1: Memory chip architecture” under Answers at the end.}
\end{keepsection}
\end{keepsection}

\begin{keepsection}
\Needspace{5\baselineskip}
\section{Exercise 3.2: Memory modules architecture}

\textbf{Connected topic:} Memory chips and modules (§5.4) \ensuremath{\cdot} Memory interleaving (§5.5)

\begin{keepsection}
\Needspace{7\baselineskip}
\subsection{3.2.1 Module configuration}

\Needspace{4\baselineskip}
Try first: Complete the chip and address-width specifications for a word-addressable 1G \ensuremath{\times} 32-bit memory built from 256M \ensuremath{\times} 4-bit chips.

\nopagebreak[4]{\small\itshape Answers: see “Exercise 3.2: Memory modules architecture” under Answers at the end.}
\end{keepsection}

\begin{keepsection}
\Needspace{7\baselineskip}
\subsection{3.2.2 Word 14 across four modules}

\Needspace{4\baselineskip}
Try first: Locate word address 14 under high-order and low-order interleaving, numbering modules from zero.

\nopagebreak[4]{\small\itshape Answers: see “Exercise 3.2: Memory modules architecture” under Answers at the end.}
\end{keepsection}

\begin{keepsection}
\Needspace{7\baselineskip}
\subsection{3.2.3 Word 300,000,000}

\Needspace{4\baselineskip}
Try first: Find the module and local word address for word address 300,000,000 under both interleaving schemes.

\nopagebreak[4]{\small\itshape Answers: see “Exercise 3.2: Memory modules architecture” under Answers at the end.}
\end{keepsection}
\end{keepsection}

\begin{keepsection}
\Needspace{5\baselineskip}
\section{Exercise 3.3: Byte-addressable memory modules}

\textbf{Connected topic:} Addressing and capacity (§5.2) \ensuremath{\cdot} Memory interleaving (§5.5)

\begin{keepsection}
\Needspace{7\baselineskip}
\subsection{3.3.1 Full-memory address width}

\Needspace{4\baselineskip}
Try first: For the same 1G \ensuremath{\times} 32-bit memory, find the number of address bits when it is byte-addressable.

\nopagebreak[4]{\small\itshape Answers: see “Exercise 3.3: Byte-addressable memory modules” under Answers at the end.}
\end{keepsection}

\begin{keepsection}
\Needspace{7\baselineskip}
\subsection{3.3.2 Word 14 after byte-addressing}

\Needspace{4\baselineskip}
Try first: Locate word 14 under both interleaving schemes. The memory modules still store complete 32-bit words.

\nopagebreak[4]{\small\itshape Answers: see “Exercise 3.3: Byte-addressable memory modules” under Answers at the end.}
\end{keepsection}
\end{keepsection}

\begin{keepsection}
\Needspace{5\baselineskip}
\section{Exercise 3.4: Endianness bytes}

\textbf{Connected topic:} Endianness and aligned access (§5.3)

\begin{keepsection}
\Needspace{6\baselineskip}
\subsection{3.4.1 Define endianness}

\Needspace{3\baselineskip}
Try first: State the definition of endianness in one sentence.

\nopagebreak[4]{\small\itshape Answers: see “Exercise 3.4: Endianness bytes” under Answers at the end.}
\end{keepsection}

\begin{keepsection}
\Needspace{7\baselineskip}
\subsection{3.4.2 Transfer between different byte orders}

\Needspace{4\baselineskip}
Try first: Explain what must happen when data moves from little-endian memory to big-endian memory.

\nopagebreak[4]{\small\itshape Answers: see “Exercise 3.4: Endianness bytes” under Answers at the end.}
\end{keepsection}

\begin{keepsection}
\Needspace{8\baselineskip}
\subsection{3.4.3 Transfer 0x12345678 unchanged}

\Needspace{5\baselineskip}
Try first: A 32-bit integer 0x12345678 is stored little-endian and copied unchanged into big-endian memory. Find the value read there and decide whether the original can be recovered without byte-order information.

\nopagebreak[4]{\small\itshape Answers: see “Exercise 3.4: Endianness bytes” under Answers at the end.}
\end{keepsection}
\end{keepsection}


\clearpage
\section{Answers}
\begin{keepsection}\subsection{Exercise 3.1: Memory chip architecture}

\begin{workedbox}{3.1.1 Memory size in bits \ensuremath{\cdot} worked solution}
\Needspace{7\baselineskip}
\begin{enumerate}
\item In the notation 4G \ensuremath{\times} 32, there are 4G addressable words and 32 bits per word. For this exercise, G means 2\textsuperscript{30}, so 4G = 2\textsuperscript{32} words.
\item Multiply words by bits per word: 2\textsuperscript{32} \ensuremath{\times} 32 = 2\textsuperscript{32} \ensuremath{\times} 2\textsuperscript{5} = 2\textsuperscript{37} bits. Thus x = 37.
\item The check is dimensional: word count \ensuremath{\times} bits/word gives bits. The 32-bit width adds five powers of two to the 4G word count.
\end{enumerate}
\end{workedbox}

\end{keepsection}
\begin{keepsection}\subsection{Exercise 3.1: Memory chip architecture}

\begin{workedbox}{3.1.2 Word-addressable memory \ensuremath{\cdot} worked solution}
\Needspace{5\baselineskip}
\begin{enumerate}
\item The memory contains 4G = 2\textsuperscript{32} words. Word-addressing assigns one distinct address to each word.
\item The number of address bits is log\textsubscript{2}(2\textsuperscript{32}) = 32 bits.
\item A 32-bit address has 2\textsuperscript{32} patterns, exactly enough for all word addresses from 0 through 2\textsuperscript{32}\ensuremath{-}1.
\end{enumerate}
\end{workedbox}

\end{keepsection}
\begin{keepsection}\subsection{Exercise 3.1: Memory chip architecture}

\begin{workedbox}{3.1.3 Byte-addressable memory \ensuremath{\cdot} worked solution}
\Needspace{6\baselineskip}
\begin{enumerate}
\item Each 32-bit word contains 4 bytes. The number of bytes is 2\textsuperscript{32} words \ensuremath{\times} 4 bytes/word = 2\textsuperscript{34} bytes.
\item Addressing 2\textsuperscript{34} individual bytes requires 34 address bits.
\item This is two more bits than word addressing because each word contains 2\textsuperscript{2} byte locations.
\end{enumerate}
\end{workedbox}

\end{keepsection}
\begin{keepsection}\subsection{Exercise 3.1: Memory chip architecture}

\begin{workedbox}{3.1.4 Byte address range \ensuremath{\cdot} worked solution}
\Needspace{8\baselineskip}
\begin{enumerate}
\item There are 2\textsuperscript{34} byte addresses, numbered consecutively from zero. The highest address is therefore the count minus one.
\item The lowest address is 0x000000000. The highest is 2\textsuperscript{34}\ensuremath{-}1 = 0x3FFFFFFFF. The address width is 34 bits.
\item Check that the next address, 0x400000000, is 2\textsuperscript{34} and is outside the memory. Leading zeroes in the lowest address are only padding.
\end{enumerate}
\end{workedbox}

\end{keepsection}
\begin{keepsection}\subsection{Exercise 3.2: Memory modules architecture}

\begin{workedbox}{3.2.1 Module configuration \ensuremath{\cdot} worked solution}
\Needspace{12\baselineskip}
\begin{enumerate}
\item A chip supplies 4 bits per word, so eight chips in parallel supply the required 32-bit word. Each module then has 256M = 2\textsuperscript{28} words. Four such modules provide 1G = 2\textsuperscript{30} words.
\item The full system uses 8 \ensuremath{\times} 4 = 32 chips. A module needs log\textsubscript{2}(2\textsuperscript{28}) = 28 row-address bits; the full memory needs log\textsubscript{2}(2\textsuperscript{30}) = 30 word-address bits.
\item The configuration is:
\item \Needspace{10\baselineskip}
\begin{center}\small\begin{tabularx}{\linewidth}{>{\raggedright\arraybackslash}X>{\raggedright\arraybackslash}X}\toprule
\textbf{Specification} & \textbf{Quantity} \\ \midrule
RAM chips per memory module & 8 \\
Memory modules & 4 \\
RAM chips for full memory & 32 \\
Address bits per module & 28 \\
Address bits for full memory & 30 \\
\bottomrule\end{tabularx}\end{center}
\item Check capacity and width separately: 4 modules \ensuremath{\times} 256M words = 1G words; 8 chips \ensuremath{\times} 4 bits = 32 bits per word.
\end{enumerate}
\end{workedbox}

\end{keepsection}
\begin{keepsection}\subsection{Exercise 3.2: Memory modules architecture}

\begin{workedbox}{3.2.2 Word 14 across four modules \ensuremath{\cdot} worked solution}
\Needspace{9\baselineskip}
\begin{enumerate}
\item There are four modules, so two address bits select a module. Each module holds 2\textsuperscript{28} words. High-order interleaving uses the most significant address bits for the module; low-order interleaving uses the least significant two bits.
\item High-order: module = floor(14/2\textsuperscript{28}) = 0; local word address = 14. Low-order: module = 14 mod 4 = 2; local word address = floor(14/4) = 3.
\item Check by reconstructing the addresses. High-order module 0 starts at word 0, so local 14 is global 14. Low-order module 2 stores words 2, 6, 10, 14, …; its fourth row is local address 3.
\end{enumerate}
\end{workedbox}

\end{keepsection}
\begin{keepsection}\subsection{Exercise 3.2: Memory modules architecture}

\begin{workedbox}{3.2.3 Word 300,000,000 \ensuremath{\cdot} worked solution}
\Needspace{9\baselineskip}
\begin{enumerate}
\item For high-order interleaving, divide by the module capacity 2\textsuperscript{28} = 268,435,456. For low-order interleaving, divide by four modules.
\item High-order: 300,000,000 = 1 \ensuremath{\times} 268,435,456 + 31,564,544, so the location is module 1, local word address 31,564,544. Low-order: 300,000,000 mod 4 = 0 and floor(300,000,000/4) = 75,000,000, so the location is module 0, local word address 75,000,000.
\item Reconstruct: high-order 1 \ensuremath{\times} 2\textsuperscript{28} + 31,564,544 = 300,000,000; low-order 4 \ensuremath{\times} 75,000,000 + 0 = 300,000,000.
\end{enumerate}
\end{workedbox}

\end{keepsection}
\begin{keepsection}\subsection{Exercise 3.3: Byte-addressable memory modules}

\begin{workedbox}{3.3.1 Full-memory address width \ensuremath{\cdot} worked solution}
\Needspace{7\baselineskip}
\begin{enumerate}
\item The capacity is 1G words \ensuremath{\times} 32 bits/word. Each word is four bytes, so there are 1G \ensuremath{\times} 4 = 4G = 2\textsuperscript{32} bytes.
\item The full byte-addressable memory therefore needs log\textsubscript{2}(2\textsuperscript{32}) = 32 address bits.
\item Check against word-addressing: 1G words require 30 bits; the two byte-within-word bits raise the byte address to 32 bits.
\end{enumerate}
\end{workedbox}

\end{keepsection}
\begin{keepsection}\subsection{Exercise 3.3: Byte-addressable memory modules}

\begin{workedbox}{3.3.2 Word 14 after byte-addressing \ensuremath{\cdot} worked solution}
\Needspace{10\baselineskip}
\begin{enumerate}
\item Word index 14 begins at byte address 14 \ensuremath{\times} 4 = 56. Because each module stores a complete 32-bit word, select the module using the word index; the two lowest byte-address bits identify the byte within that word.
\item High-order interleaving: word 14 is in module 0 at local word row 14, beginning at local byte address 56. Low-order interleaving: module = 14 mod 4 = 2 and local word row = floor(14/4) = 3; the word begins at local byte address 3 \ensuremath{\times} 4 = 12. Its global byte address is still 56.
\item The low-order module bits come from the word address after removing the two byte-offset bits. Reconstruct the global word index: 4 \ensuremath{\times} 3 + 2 = 14. The first byte of that word is 56.
\end{enumerate}
\end{workedbox}

\end{keepsection}
\begin{keepsection}\subsection{Exercise 3.4: Endianness bytes}

\begin{workedbox}{3.4.1 Define endianness \ensuremath{\cdot} worked solution}
\Needspace{8\baselineskip}
\begin{enumerate}
\item A multi-byte value occupies several byte addresses, so the system needs a rule for which byte comes first.
\item Endianness is the order in which a multi-byte value’s bytes are stored in memory: little-endian places the least significant byte at the lowest address, while big-endian places the most significant byte there.
\item This rule affects multi-byte values; a single-byte value has no byte-order choice.
\end{enumerate}
\end{workedbox}

\end{keepsection}
\begin{keepsection}\subsection{Exercise 3.4: Endianness bytes}

\begin{workedbox}{3.4.2 Transfer between different byte orders \ensuremath{\cdot} worked solution}
\Needspace{9\baselineskip}
\begin{enumerate}
\item A raw byte copy preserves the sequence of bytes. The receiving system then interprets that sequence using its own byte-order rule.
\item To preserve the numeric value, the sender and receiver must agree on a serialized byte order, or one side must reorder the bytes during transfer. If bytes are copied unchanged without conversion, a multi-byte number may be read as a different value.
\item Single-byte text is unaffected. A robust file or network format specifies byte order rather than relying on the sender’s native layout.
\end{enumerate}
\end{workedbox}

\end{keepsection}
\begin{keepsection}\subsection{Exercise 3.4: Endianness bytes}

\begin{workedbox}{3.4.3 Transfer 0x12345678 unchanged \ensuremath{\cdot} worked solution}
\Needspace{8\baselineskip}
\begin{enumerate}
\item In little-endian order, the bytes at increasing addresses are 78, 56, 34, 12. The transfer keeps that sequence unchanged.
\item A big-endian receiver treats the first byte as most significant, so it reads 0x78563412.
\item The original value cannot be determined from those bytes alone without knowing the sender’s byte order or other format metadata: the same four bytes can represent different integers under different conventions.
\end{enumerate}
\end{workedbox}

\end{keepsection}

\vfill{\small Personal worked solutions based on the APDR Blackboard Exercise 3 sheet, written in my own words; no exercise-sheet prose or figures are reproduced. Not an official FAU publication; check course details against the current materials. Made by \textbf{Thanate Amornjiraporn (Focus)}.}

\end{document}
